\section{Experimental design}
\label{sec:design}

\subsection{Data}
\begin{table}[tbp]
\centering\scriptsize
\caption{Real response matrices. Verdict: one deterministic evaluation per system and item (single) or a recorded pass rate over several samples binarized at one half (aggregate; TruthfulQA: continuous MC2 score binarized). $J$ systems, $F$ heuristic families, $I$ items; accuracy range across systems; items solved by no system. Role: confirmatory in studies 1 and 2, or development bank used only to design the protocols (Sections~\ref{sec:dev} and~\ref{sec:design2}). SE: software engineering.}
\label{tab:data}
\setlength{\tabcolsep}{2.6pt}
\begin{tabular}{@{}lllllrrrcr@{}}
\toprule
Bank & Role & Domain & Oracle & Verdict & $J$ & $F$ & $I$ & Accuracy & Unsolved \\
\midrule
GSM8K & 1, 2 & math & numeric match & single & 395 & 227 & 1319 & 0.00--0.77 & 13 \\
ARC-C & 1, 2 & multiple choice & log-lik.\ choice & single & 395 & 227 & 1172 & 0.26--0.73 & 51 \\
HellaSwag & 1, 2 & multiple choice & log-lik.\ choice & single & 395 & 227 & 2000 & 0.31--0.90 & 43 \\
WinoGrande & 1, 2 & multiple choice & log-lik.\ choice & single & 395 & 227 & 1267 & 0.59--0.87 & 1 \\
MMLU & 1, 2 & multiple choice & log-lik.\ choice & single & 395 & 227 & 2000 & 0.26--0.80 & 0 \\
SWE-bench V. & 1, 2 & agentic SE & repository tests & single & 175 & 123 & 500 & 0.00--0.79 & 26 \\
LiveCodeBench & 1, 2 & code generation & hidden tests & aggregate & 28 & 15 & 1055 & 0.23--0.87 & 46 \\
CRUXEval-O & 1, 2 & code reasoning & execution & aggregate & 31 & 12 & 800 & 0.22--0.78 & 98 \\
SWE-bench Lite & dev. & agentic SE & repository tests & single & 84 & 49 & 300 & 0.00--0.60 & 35 \\
CRUXEval-I & dev. & code reasoning & execution & aggregate & 31 & 12 & 800 & 0.14--0.78 & 59 \\
TruthfulQA & dev. & multiple choice & log-lik.\ choice & aggregate & 395 & 227 & 817 & 0.30--0.73 & 47 \\
\bottomrule
\end{tabular}
\end{table}

We use public item-level verdicts of real systems (Table~\ref{tab:data}). Open LLM
Leaderboard \citep{beeching2023openllm} verdicts for 395 models on GSM8K
\citep{cobbe2021gsm8k}, ARC-Challenge \citep{clark2018arc}, HellaSwag
\citep{zellers2019hellaswag}, WinoGrande \citep{sakaguchi2020winogrande}, MMLU
\citep{hendrycks2021mmlu} and TruthfulQA \citep{lin2022truthfulqa} are taken from the
release of \citet{polo2024tinybenchmarks}; HellaSwag and MMLU are uniformly subsampled to
2,000 items with a fixed seed, and TruthfulQA's continuous MC2 score is binarized at 0.5.
SWE-bench Verified and Lite \citep{jimenez2024swebench} verdicts (instance resolved by the
official harness) come from the public submissions repository (175 and 84 systems).
LiveCodeBench \citep{jain2025livecodebench} pass rates for 28 models on 1,055 problems and
CRUXEval \citep{gu2024cruxeval} pass rates of temperature-0.2 samples for 31 models are
binarized at one half; these two banks therefore carry recorded aggregate verdicts
(Section~\ref{sec:framework}), the other six deterministic ones. Every source is pinned to
a git commit and every processed matrix carries a SHA-256 hash.

\paragraph{Synthetic banks.} Three simulated banks with known abilities complement the
real data (300 systems, 600 items each, $\theta\sim\mathcal N(0,1)$,
$\log a\sim\mathcal N(0,0.5^2)$, $b\sim\mathcal N(0,1.5^2)$): a well-specified 2PL bank; a
bank in which 60\% of the items have a guessing floor of 0.25; and a bank in which half of
the items depend on a second ability dimension correlated 0.6 with the first. The fitted
model is always the unidimensional 2PL. Synthetic results never enter a confirmatory
decision.

\paragraph{Families and strict clusters.} Systems that share an underlying model are
not independent. \emph{Heuristic families} are defined by organization namespace (Open
LLM Leaderboard), by model line including reasoning-effort and version variants
(LiveCodeBench, CRUXEval), and, for SWE-bench, by the set of underlying LLMs declared in a
submission's metadata (two agentic systems are related if they share any LLM). Heuristic
families miss renamed copies, merges and fine-tunes across organizations, so study~2 uses
\emph{strict clusters}: connected components of the graph that joins two systems if they
share a heuristic family or if the correlation of their standardized 2PL residuals is at
least 0.5---they share at least half of their idiosyncratic error pattern beyond what
ability explains. Residuals are pooled over banks that share systems (the six Open LLM
Leaderboard banks; CRUXEval-O and -I). Agreement measures such as Cohen's $\kappa$ were
not usable for this purpose: on HellaSwag, $\kappa\ge0.8$ chains 353 of 395 systems into
one component because all strong models agree on most items. The strict clusters number
\ClusOLLB\ on the Open LLM Leaderboard (largest \ClusOLLBMax\ systems), \ClusSWE\ on
SWE-bench Verified, \ClusLCB\ on LiveCodeBench and \ClusCRUX\ on CRUXEval.

\subsection{Study 1: frozen plan, leave-family-out replay}\label{sec:dev}
For each bank, up to 60 targets were drawn stratified by accuracy decile (all systems
for LiveCodeBench and CRUXEval). For each target and replicate $r=1,\dots,5$, a holdout of
20\% of the items was drawn and shared by all policies and panel sizes; for each panel
size $N\in\{4,8,16,32,64\}$ a panel was drawn outside the target's heuristic family. The
twelve policies of Section~\ref{sec:framework} (five criteria under both models, PWFI
under the marginal model only, plus Naive) were replayed on the target's recorded
verdicts with a budget of $T=100$ items. All random streams derive from one master seed
by hashing (bank, purpose, target, replicate, $N$, policy).

The protocol was designed on development banks only (SWE-bench Lite, CRUXEval-I,
TruthfulQA, synthetic 2PL; Section~\ref{sec:devhist}), where one problem appeared: with
fixed weakly-informative priors, items solved by no calibrator were shrunk towards an
arbitrary difficulty, so empirical-Bayes hyperpriors were adopted for every policy. The
plan (banks, protocol, endpoints, four hypotheses, tests, multiplicity control, decision
rule), the analysis code and the data manifests were then hashed and committed before any
confirmatory run, and the run wrapper refuses to start if a frozen file changes. No public
registry or external time stamp was used, so we call this a \emph{frozen analysis plan}
rather than a preregistration \citep{nosek2018preregistration}; the git history shows the
freeze commit before every confirmatory artifact.

\paragraph{Hypotheses} (primary panel size $N=8$):
\textbf{H1} EIG:mg has lower integrated holdout log loss than Entropy:mg;
\textbf{H2} EIG:mg has lower integrated holdout log loss than MFI:pi;
\textbf{H3} EIG:mg has lower integrated score error than Naive subsampling;
\textbf{H4} EIG:mg's 90\% score intervals are better calibrated than MFI:pi's
($|\mathrm{coverage}-0.9|$). Each contrast is a two-sided Wilcoxon signed-rank test
\citep{wilcoxon1945individual} on paired per-target values (replicates averaged within
target), with Holm correction \citep{holm1979simple} over the whole family
(8 banks $\times$ panel sizes $\times$ 4 hypotheses); a hypothesis is \emph{supported} if,
at $N=8$, the adjusted test favours EOST-M in at least five of the eight banks and the
baseline in none. Because targets of one family are not independent, every study-1
contrast is also re-analysed post hoc with heuristic families and with strict clusters as
units (cluster means, the same tests and decision rule).

\subsection{Study 2: independent targets and strong competitors}\label{sec:design2}
Study~2 was designed after an external review of study~1 and frozen before execution
(plan, hypotheses, cluster file and code hashed and committed; the plan and its hash were
also sent to the first author, outside the execution environment, before the first run).
It differs from
study~1 in four ways. \emph{Targets} are one randomly chosen representative per strict
cluster (systems that were not study-1 targets preferred), at most 40 per bank, stratified
by accuracy, so targets are independent units; \emph{panels} ($N\in\{4,8,16\}$, three
replicates, new seeds) exclude the target's whole strict cluster. \emph{Competitors} are
the strongest we could identify:
\begin{itemize}\itemsep0pt
\item \textsc{MFI:1pr}, a regularized Rasch-type plug-in CAT (one common discrimination
with $\log a\sim\mathcal N(0,0.05^2)$, empirical-Bayes prior on $b$), and
\textsc{EIG:1prmg}, the same model integrated (both found post hoc in study~1);
\item \textsc{FB:bmi} and \textsc{FB:eig}, a \emph{posterior-draw marginal Bayesian CAT} (PDM
Bayesian CAT; called the fully Bayesian CAT in the frozen plan): a Gibbs sampler for the
hierarchical 2PL of the panel (abilities, item parameters on the grid, hyperparameters,
with a translation move for the overall location), eight retained draws, and the
target's posterior computed with each item's response function averaged over the draws
separately (the target of \citealp{vanderlinden2020fast}; Lemma~\ref{lem:exact} shows that
this integrates the calibrators' abilities and the hyperprior item by item, not jointly);
selection by Bayesian maximum information
(posterior-weighted expected Fisher information, the strategy recommended by
\citealp{fink2025accounting}) or by EIG ($N\in\{4,8\}$);
\item \textsc{Random:mg}, random items scored with the EOST-M model, which isolates the
value of adaptivity from that of the calibration;
\item \textsc{ATLAS}, an ATLAS-style CAT \citep{li2026atlas}: 3PL with the lower asymptote
fixed at the chance rate of the item format (ARC-C, HellaSwag, MMLU $0.25$, WinoGrande
$0.5$, others 0), first item closest in difficulty to 0, then a uniform draw among the five
most informative items at the expected a posteriori (EAP) ability \citep{kingsbury1989procedures};
\item \textsc{tinyBenchmarks} \citep{polo2024tinybenchmarks}: anchor items by K-Means on
the panel's item embeddings and the gp-IRT estimator, at budgets of 25, 50 and 100 items;
\textsc{MINCE} \citep{das2026mince}: Monte Carlo sizing on the panel and the plain mean of
one frozen random subset, at the same budgets and at its own size $n^\star$;
\item \textsc{MFI:pi} and \textsc{Naive} as in study~1, two empirical-Bayes variants of
EOST-M (cross-fitted hyperprior; Student-$t$ hyperprior family) and an exploratory mixed
policy (every fourth item uniform at random).
\end{itemize}
\emph{Intervals} of EOST-M and of the Rasch-type test are also reported after the
leave-one-calibrator-out recalibration of Section~\ref{sec:framework}. \emph{Hypotheses}
(all at $N=8$, Holm over 11 superiority hypotheses $\times$ 8 banks; supported with
$\ge5$ banks for EOST-M and none against): replications of H2--H4 with independent targets
(R2--R4); EOST-M versus the Rasch-type plug-in (C1) and marginal (C2) tests on log loss;
versus Random:mg on score error (C4); recalibrated versus raw EOST-M intervals (C5b) and
versus Naive's Wilson intervals (C6) on the interval score; versus ATLAS-style CAT on score
error (C7); versus tinyBenchmarks (C8) and MINCE (C9) on the absolute score error after 50
items. In addition, the PDM Bayesian CAT and EOST-M are declared practically equivalent
(C3) if the 90\% bootstrap interval of their relative log-loss difference lies within
$\pm1\%$ on at least six banks, and the recalibrated EOST-M interval is declared calibrated
(C5a) if its mean coverage lies in $[0.85,0.95]$ on at least six banks. A \emph{temporal
component} on SWE-bench Verified uses the latest quarter of submissions by date as targets
and only earlier submissions as calibrators ($N=8$); the synthetic banks are a secondary
component. During the run, before any result was inspected, an independent review of the
code found two problems---tinyBenchmarks can use fewer anchors than its budget, and the
burn-in of the Bayesian sampler was too short on some banks---whose post hoc
corrections were specified in the research log before the results were looked at; the
frozen run was completed unchanged and every frozen decision is reported as computed
(Section~\ref{sec:deviations}, deviation xiii).

\paragraph{Implementation of the competitors.} The \emph{Bayesian CATs} run a Gibbs sampler on the
panel's hierarchical 2PL: (1) every item's grid point is drawn from its exact conditional
grid posterior; (2) every calibrator's ability is drawn by griddy Gibbs on 65 points
spanning $\pm0.8$ around its current value; (2b) all abilities and difficulties and the
difficulty mean are shifted by a common multiple of the difficulty grid step (at most 16
steps either way) drawn from its conditional, which only involves the ability prior (the
likelihood and the difficulty hyperprior are invariant); shifts that would move any
sampled difficulty off the grid are excluded, so the move is exact only within these
bounds and rarely acts when many difficulties sit near the edge of the grid; (3) the
hyperparameters are drawn from their conjugate conditionals (flat means, Jeffreys standard
deviations). The chain starts at EOST-M's estimates; in the frozen plan, 8 draws are kept
every 4 sweeps after a burn-in of 10 sweeps, and each contributes its exact conditional item
posterior (Rao--Blackwellization). As noted above, this burn-in proved too short on some banks; the
post hoc correction (60 sweeps) is described in Section~\ref{sec:s2fb}. In 300-sweep chains on the
development banks (\texttt{scripts/fb\_mixing.py}), the translation move removed the slow
mixing of the overall location: the lag-3 autocorrelation of the mean calibrator ability
fell from \MixLocWithoutMin--\MixLocWithoutMax\ without the move to at most \MixLocWithMax\
with it. That location is genuinely uncertain---its posterior standard deviation was
\MixLocSdFour, \MixLocSdEight\ and \MixLocSdSixteen\ for $N=4$, 8 and 16, close to
$1/\sqrt N$ as implied by the $\mathcal N(0,1)$ ability prior---but it shifts abilities and
difficulties together. Scale hyperparameters mix more slowly (lag-3 autocorrelation of
$\sigma_a$ between \MixScaleMin\ and \MixScaleMax), so the eight retained draws represent
their uncertainty only roughly. Follow-up RF2 (Section~\ref{sec:rf2}) measured this with
four long chains, and exploratory follow-ups added the missing moves.

The \emph{ATLAS-style} test fits a 3PL with fixed lower asymptote by grid MAP given the
panel's abilities and the empirical-Bayes prior. \emph{tinyBenchmarks} uses K-Means (best of
five k-means++ restarts) on $(a_i,a_ib_i)$ of the panel's plug-in 2PL fit, the item nearest
to each centroid as anchor, cluster fractions as weights, the bounded maximum-likelihood
ability for p-IRT, and $\lambda=\hat b^2/(\hat v/(4n)+\hat b^2)$ with $\hat b$ the
leave-one-calibrator-out error of the IRT prediction and $\hat v$ the mean Bernoulli
variance of the calibrators. \emph{MINCE} draws 2,000 random subsets per candidate size
(steps of 25) and stops when the 95th percentile of the worst-case drift over calibrators
improves by less than one point.

\subsection{From the pilot program to the frozen plans}\label{sec:devhist}\label{sec:pilot}

\paragraph{Pilot program (historical).} Before study~1, the framework (EOST in our project) was developed through a staged pilot
program on math word problems with \texttt{Qwen3-8B} as target, greedy decoding
and exact numeric oracles, including self-generated items with an independent oracle
(the self-testing aspect of the framework). Every stage had an explicit pass/fail gate
fixed in advance. We report the stages \emph{as recorded in the frozen project log},
including failed gates and the subsequent corrections, because they motivated
the design of Sections~\ref{sec:framework}--\ref{sec:design}. Stages S0--S4 were
executed before the project's Docker-only rule; the S4B run of one model was executed in
the project's earlier container runtime, which differs from the runtime of this paper
(Section~\ref{sec:repro}). Their raw artifacts are archived with SHA-256 hashes but they
were \emph{not} re-executed for this paper, and none of the paper's claims rests on them.
Table~\ref{tab:pilot} summarizes them.

\begin{table}[htbp]
\centering\footnotesize
\caption{Pilot program (historical values from the frozen project log; not re-executed here).
AUIC: area under the information curve; IG: expected-information-gain policy (EOST).}
\label{tab:pilot}
\setlength{\tabcolsep}{4pt}
\begin{tabular}{@{}p{0.9cm}p{4.2cm}p{9.4cm}@{}}
\toprule
Stage & Setting & Recorded outcome and gate \\
\midrule
S0 & Synthetic latent-trait simulator, 30 seeds; IG vs.\ entropy, random, difficulty-matched &
IG $>$ entropy on entropy reduction (6.77 vs.\ 6.51; Holm $p=2.8\times10^{-9}$) and AUIC (5.29 vs.\ 5.12; rank-biserial 1.0). \textbf{Pass}. \\
S1-R2 & GSM-Symbolic, 45 items; 4 calibration conditions (two Qwen3 checkpoints $\times$ thinking on/off) &
Conditions were not independent systems. \textbf{Gate fail}. \\
S1-R3 & Same bank, regularized 2PL & IG AUIC 0.950 vs.\ entropy 0.923 (29/30 seeds); identical final holdout. Exploratory. \\
S2 & LLM-generated items, procedural oracle (75 items, 5 families) &
After fixing a numeric-tolerance bug: IG AUIC 1.224 vs.\ 1.139 (27/30 seeds, $p=1.2\times10^{-4}$); gains early in the test, endpoint saturated, log-loss tail risk. Publication gate \textbf{fail}. \\
S3 & Circularity: self-, external- and procedurally-generated items (45 matched triplets) &
Self minus external accuracy 0.089 (95\% CI $[0,0.2]$, McNemar $p=0.22$): no \emph{detectable} circularity gap at this power (not equivalence). \\
S4 & GSM8K, GSM-Hard, SVAMP, ASDiv (240 items); 4 calibration conditions &
Macro AUIC IG 0.967 vs.\ entropy 0.858 (30/30 seeds). Original gate \textbf{fail} (predictive non-inferiority). An audit found the original Brier holdout was policy-dependent; with a common fixed holdout, integrated Brier 0.0936 vs.\ 0.0969 (24/30 seeds, $p=3.1\times10^{-5}$), final Brier tie (16/30). GSM-Hard went the other way (Brier 0.195 vs.\ 0.183). 2PL weakly identified: 7 distinct $b$, 50.8\% of $a$ at the upper bound. \\
S4B & Expanding the calibration panel to independent systems &
Six calibrators passed a post-hoc parser audit (parse failures within the project's internal 5\% threshold); a seventh, community checkpoint produced 235/240 truncated generations---a generation-protocol failure, not a capability measurement---and is classified as invalid for calibration. Incomplete; not analysed. \\
\bottomrule
\end{tabular}
\end{table}

Three lessons from the pilot shaped EOST v2. First, the information advantage
of IG over predictive entropy was consistent across synthetic, self-generated
and multi-benchmark settings, but it was concentrated early in the test, as
Proposition~\ref{prop:local} predicts. Second, predictive evaluation must use a
holdout shared by all policies; the policy-dependent holdout of the original S4
analysis produced a spurious non-inferiority failure, which we keep reported as
such. Third, and most important, a calibration panel of four conditions drawn
from two checkpoints cannot identify a 2PL: half of the discriminations
saturated at the optimizer bound, the textbook setting for capitalization on
chance (Proposition~\ref{prop:cap}). The confirmatory study therefore targets
exactly this small-panel regime, with independent systems, leave-family-out
panels and calibration-marginal selection.

\paragraph{Development runs.} The development runs on the development banks that shaped both
protocols are summarized in Table~\ref{tab:dev}; the switch to empirical-Bayes hyperpriors
described in Section~\ref{sec:dev} was decided there, before the plan of study~1 was frozen.

\begin{table}[tbp]
\centering\scriptsize
\caption{Development phase (exploratory; 30 targets, 3 replicates). The switch from fixed to empirical-Bayes (EB) item hyperpriors was decided here, before freezing the confirmatory plan.}
\label{tab:dev}
\begin{tabular}{@{}llrrrrrrr@{}}
\toprule
 & & & \multicolumn{2}{c}{holdout log loss} & \multicolumn{2}{c}{score IAE} & \multicolumn{2}{c}{coverage} \\
\cmidrule(lr){4-5}\cmidrule(lr){6-7}\cmidrule(lr){8-9}
Bank & prior & $N$ & EIG:mg & MFI:pi & EIG:mg & MFI:pi & EIG:mg & MFI:pi \\
\midrule
Synth. 2PL & fixed & 4 & 0.5762 & 0.6410 & 0.0433 & 0.0543 & 0.78 & 0.55 \\
SWE-bench Lite & fixed & 4 & 0.4133 & 0.4110 & 0.0430 & 0.0406 & 0.74 & 0.69 \\
CRUXEval-I & fixed & 4 & 0.4833 & 0.5054 & 0.0342 & 0.0441 & 0.88 & 0.63 \\
TruthfulQA & fixed & 4 & 0.4215 & 0.4179 & 0.0335 & 0.0395 & 0.87 & 0.66 \\
Synth. 2PL & fixed & 8 & 0.5371 & 0.5696 & 0.0382 & 0.0427 & 0.77 & 0.59 \\
SWE-bench Lite & fixed & 8 & 0.3743 & 0.3773 & 0.0349 & 0.0292 & 0.74 & 0.75 \\
CRUXEval-I & fixed & 8 & 0.4460 & 0.4659 & 0.0315 & 0.0400 & 0.80 & 0.55 \\
TruthfulQA & fixed & 8 & 0.3726 & 0.3755 & 0.0264 & 0.0305 & 0.83 & 0.64 \\
Synth. 2PL & fixed & 64 & 0.4914 & 0.4931 & 0.0290 & 0.0275 & 0.77 & 0.80 \\
SWE-bench Lite & fixed & 64 & 0.3254 & 0.3267 & 0.0264 & 0.0250 & 0.80 & 0.81 \\
TruthfulQA & fixed & 64 & 0.3204 & 0.3214 & 0.0223 & 0.0237 & 0.70 & 0.64 \\
Synth. 2PL & EB & 4 & 0.5711 & 0.5884 & 0.0424 & 0.0475 & 0.87 & 0.78 \\
SWE-bench Lite & EB & 4 & 0.4036 & 0.4207 & 0.0393 & 0.0371 & 0.77 & 0.69 \\
CRUXEval-I & EB & 4 & 0.4833 & 0.5113 & 0.0365 & 0.0452 & 0.83 & 0.64 \\
TruthfulQA & EB & 4 & 0.4110 & 0.4323 & 0.0367 & 0.0392 & 0.77 & 0.64 \\
Synth. 2PL & EB & 8 & 0.5332 & 0.5403 & 0.0359 & 0.0407 & 0.92 & 0.82 \\
SWE-bench Lite & EB & 8 & 0.3649 & 0.3730 & 0.0296 & 0.0293 & 0.83 & 0.79 \\
CRUXEval-I & EB & 8 & 0.4450 & 0.4702 & 0.0326 & 0.0416 & 0.79 & 0.53 \\
TruthfulQA & EB & 8 & 0.3679 & 0.3851 & 0.0284 & 0.0308 & 0.78 & 0.65 \\
Synth. 2PL & EB & 64 & 0.4902 & 0.4915 & 0.0261 & 0.0273 & 0.89 & 0.87 \\
SWE-bench Lite & EB & 64 & 0.3228 & 0.3243 & 0.0237 & 0.0240 & 0.86 & 0.84 \\
TruthfulQA & EB & 64 & 0.3202 & 0.3223 & 0.0236 & 0.0248 & 0.67 & 0.63 \\
\bottomrule
\end{tabular}
\end{table}


\paragraph{Choosing the recalibration rule.} Three leave-one-calibrator-out rules were
implemented: multiplying the model half-width by the 90\% quantile of the pseudo-targets'
absolute standardized errors divided by $z_{0.95}$ (Rq) or by their root mean square
(Rr), and replacing it by the 90\% quantile of their absolute errors (Ra). On the
development banks only, and by a selection rule written in the research log before that
run, the rule whose mean coverage at $N\in\{4,8\}$ was closest to 0.90 was chosen: Ra
(\DevRaCov, against \DevRqCov\ for Rq and \DevRrCov\ for Rr; Table~\ref{tab:devrule}). The same
run showed that no rule helps with four calibrators, which is why the recalibration
hypotheses concern $N=8$.

\IfFileExists{tables/devrule.tex}{\begin{table}[tbp]
\centering\footnotesize
\caption{Choice of the interval recalibration on the development banks (before study 2; 20 targets or fewer, two replicates). Mean coverage of EOST-M's 90\% score interval without recalibration and with the three leave-one-calibrator-out rules: quantile of standardized errors (Rq), their root mean square (Rr), quantile of absolute errors (Ra). The rule written in advance selects the rule with mean coverage closest to 0.90: Ra.}
\label{tab:devrule}
\begin{tabular}{@{}lrcccc@{}}
\toprule
Bank & $N$ & raw & Rq & Rr & Ra \\
\midrule
CRUXEval-I & 4 & 0.85 & 0.81 & 0.82 & 0.84 \\
CRUXEval-I & 8 & 0.84 & 0.88 & 0.89 & 0.92 \\
TruthfulQA & 4 & 0.77 & 0.75 & 0.77 & 0.77 \\
TruthfulQA & 8 & 0.70 & 0.75 & 0.75 & 0.82 \\
SWE-bench Lite & 4 & 0.80 & 0.74 & 0.74 & 0.77 \\
SWE-bench Lite & 8 & 0.76 & 0.82 & 0.82 & 0.87 \\
Synth. 2PL & 4 & 0.85 & 0.77 & 0.80 & 0.79 \\
Synth. 2PL & 8 & 0.86 & 0.85 & 0.85 & 0.85 \\
\midrule
mean & & 0.80 & 0.80 & 0.80 & 0.83 \\
\bottomrule
\end{tabular}
\end{table}
}{}

\subsection{Deviations from the frozen plans}\label{sec:deviations}
\DeviationText

\subsection{Computation, reproducibility and environment}\label{sec:repro}
All computations of this paper are executed by \texttt{./run} inside one pinned CPU-only
container image; scientific containers have no network access, and data were fetched
beforehand with sparse git checkouts of pinned commits. Study~1 (\ConfirmRuns\ adaptive
tests of 100 items) required about \ConfirmHours~h on a two-core machine and study~2 about
\StudyTwoHours~h, plus \StudyTwoFixHours~h for its post hoc corrections.
\begin{itemize}\itemsep1pt\sloppy
\item \emph{Runtime image} \texttt{eost-runtime:cpu}, image ID \texttt{sha256:\ImageIdShort\dots}
(full ID in \texttt{configs/ENVIRONMENT.json}), built on 2026-10-29 from \texttt{ubuntu:24.04}
pinned by digest \texttt{sha256:\BaseDigestShort\dots}; Python 3.12.3, NumPy 2.5.4, SciPy
1.18.1, pandas 3.1.0; complete package list in \texttt{configs/ENVIRONMENT.json} and
\texttt{requirements.lock}. Study~1, its post hoc analyses, study~2 and the post hoc
corrections of study~2 all ran in this image; the post hoc recomputation of study~1 one week
after the frozen runs reproduced the stored rows bit for bit, and the correction runs of
study~2 reproduce the frozen Bayesian-CAT rows exactly when given the frozen burn-in.
\item \emph{Paper image}: TeX Live, image ID \texttt{sha256:\PaperImageIdShort\dots}.
\item \emph{Pilot image} (historical, Section~\ref{sec:pilot}): an earlier runtime of the
project that shares the tag \texttt{eost-runtime:cpu} but is a different image
(\texttt{python:3.12-slim-bookworm}, Python 3.12.15, PyTorch 2.7.0+cpu, image ID
\texttt{sha256:\PilotImageIdShort\dots}, on a different workstation); no result of this paper
was computed in it.
\item \emph{Sources} pinned to git commits: SWE-bench experiments \texttt{40f164d},
tinyBenchmarks \texttt{e9a8b10}, LiveCodeBench site \texttt{f186ab5}, CRUXEval
\texttt{190faf1}, and, for the live protocol only, GSM8K \texttt{3101c7d}. Processed
matrices, both frozen plans and every artifact are SHA-256 hashed.
\item \emph{Tests and checks}: unit tests cover the analytic gradients, parameter recovery,
the empirical-Bayes estimator, the information cap for the logistic link
(Proposition~\ref{prop:cap}), the small-variance limit of EIG, the Jensen inequality of
Corollary~\ref{cor:jensen}, the EVR closed form and its numerical guard, the equivalence of
the study-2 replay and live engines with the frozen replay, the Gibbs sampler, its post hoc
copy and the reduction baselines. Every number in the text is a macro generated from the
artifacts, and every qualitative statement about results in
Sections~\ref{sec:study1}--\ref{sec:conclusion} is encoded as a predicate that
\texttt{scripts/check\_claims.py} re-evaluates on the artifacts
(\texttt{CLAIMS\_GATE=PASS}; \texttt{artifacts/CLAIMS.json}).
\end{itemize}
